\documentclass[aps,preprint]{revtex4}%
\usepackage{amsfonts}
\usepackage{amsmath}
\usepackage{amssymb}
\usepackage{graphicx}%
\setcounter{MaxMatrixCols}{30}
%TCIDATA{OutputFilter=latex2.dll}
%TCIDATA{Version=4.00.0.2321}
%TCIDATA{CSTFile=revtex4.cst}
%TCIDATA{Created=Thursday, February 02, 2006 16:37:04}
%TCIDATA{LastRevised=Thursday, February 02, 2006 23:52:50}
%TCIDATA{<META NAME="GraphicsSave" CONTENT="32">}
%TCIDATA{<META NAME="DocumentShell" CONTENT="Articles\SW\REVTeX 4">}
\newtheorem{theorem}{Theorem}
\newtheorem{acknowledgement}[theorem]{Acknowledgement}
\newtheorem{algorithm}[theorem]{Algorithm}
\newtheorem{axiom}[theorem]{Axiom}
\newtheorem{claim}[theorem]{Claim}
\newtheorem{conclusion}[theorem]{Conclusion}
\newtheorem{condition}[theorem]{Condition}
\newtheorem{conjecture}[theorem]{Conjecture}
\newtheorem{corollary}[theorem]{Corollary}
\newtheorem{criterion}[theorem]{Criterion}
\newtheorem{definition}[theorem]{Definition}
\newtheorem{example}[theorem]{Example}
\newtheorem{exercise}[theorem]{Exercise}
\newtheorem{lemma}[theorem]{Lemma}
\newtheorem{notation}[theorem]{Notation}
\newtheorem{problem}[theorem]{Problem}
\newtheorem{proposition}[theorem]{Proposition}
\newtheorem{remark}[theorem]{Remark}
\newtheorem{solution}[theorem]{Solution}
\newtheorem{summary}[theorem]{Summary}
\newenvironment{proof}[1][Proof]{\noindent\textbf{#1.} }{\ \rule{0.5em}{0.5em}}
\begin{document}
\preprint{ }
\title[ ]{Partitioning }
\author{Brian Weiner}
\affiliation{Penn State Univ}
\startpage{1}
\maketitle
\tableofcontents


\subsection{General Form of Partitioning}

Consider the matrix equation
\begin{equation}
\left(
\begin{array}
[c]{cc}%
\mathbb{A}_{11} & \mathbb{A}_{12}\\
\mathbb{A}_{21} & \mathbb{A}_{22}%
\end{array}
\right)  \left(
\begin{array}
[c]{c}%
\mathbf{x}_{1}\\
\mathbf{x}_{2}%
\end{array}
\right)  =\left(
\begin{array}
[c]{c}%
\mathbf{y}_{1}\\
\mathbf{y}_{2}%
\end{array}
\right)  \label{ME1}%
\end{equation}
that is equivalent to the coupled equations
\begin{align}
\mathbb{A}_{11}\mathbf{x}_{1}+\mathbb{A}_{12}\mathbf{x}_{2}  &  =\mathbf{y}%
_{1}\nonumber\\
\mathbb{A}_{21}\mathbf{x}_{1}+\mathbb{A}_{22}\mathbf{x}_{2}  &  =\mathbf{y}%
_{2}%
\end{align}
we can solve for $\mathbf{x}_{1}$ from the top equation to obtain a ''fixed''
$\mathbf{x}_{1}$
\begin{equation}
\mathbf{x}_{1}=\mathbb{A}_{11}^{-1}\left(  \mathbf{y}_{1}-\mathbb{A}%
_{12}\mathbf{x}_{2}\right)  \label{x1eq}%
\end{equation}
and then substitute in the bottom equation to obtain a solution for
$\mathbf{x}_{2}$
\begin{equation}
\mathbf{x}_{2}=\left(  \mathbb{A}_{22}+\mathbb{A}_{21}\mathbb{A}_{11}%
^{-1}\mathbb{A}_{12}\right)  ^{-1}\left(  \mathbf{y}_{2}-\mathbb{A}%
_{21}\mathbb{A}_{11}^{-1}\mathbf{y}_{1}\right)  \label{x2sol}%
\end{equation}
which then determines $\mathbf{x}_{1}$ by eq. $\left(  \ref{x1eq}\right)  $.
The choice of which variables to call $\mathbf{x}_{1}$ and $\mathbf{x}_{2}$ in
any problem is of course quite arbitrary. One should also observe that the
$\mathbf{x}_{2}$ obtained from eq. $\left(  \ref{x2sol}\right)  $ can also be
expressed as
\begin{equation}
\mathbf{x}_{2}=\mathbb{A}_{22}^{-1}\left(  \mathbf{y}_{2}-\mathbb{A}%
_{21}\mathbf{x}_{1}\right)  \label{x2leq}%
\end{equation}
where $\mathbf{x}_{1}$ is given by eq. $\left(  \ref{x1eq}\right)  $. If
neither $\mathbb{A}$ nor $\mathbf{y}$ depend on $\mathbf{x}$ and all the
inverses exist then a fixed solution for $\mathbf{x}_{1}$ and $\mathbf{x}_{2}$
is obtained in one iteration.

The equations we deal with are, however, in general nonlinear i.e.
\begin{align}
\mathbb{A}_{ij}  &  =\mathbb{A}_{ij}\left(  \mathbf{x}_{1},\mathbf{x}%
_{2}\right)  ;\quad1\leq i,j\leq2\nonumber\\
\mathbf{y}_{i}  &  =\mathbf{y}_{i}\left(  \mathbf{x}_{1},\mathbf{x}%
_{2}\right)  ;\quad1\leq i\leq2
\end{align}
and we have to iterate to obtain a solution by fixing a value for
$\mathbf{x}_{1}$ and solving for $\mathbf{x}_{2\text{ }}$to obtain
\begin{equation}
\mathbf{x}_{2}=\left(  \mathbb{A}_{22}+\mathbb{A}_{21}\mathbb{A}_{11}%
^{-1}\mathbb{A}_{12}\right)  ^{-1}\left(  \mathbf{y}_{2}-\mathbb{A}%
_{21}\mathbb{A}_{11}^{-1}\mathbf{y}_{1}\right)  =\mathbb{A}_{22}^{-1}\left(
\mathbf{y}_{2}-\mathbb{A}_{21}\mathbf{x}_{1}\right)  \label{x2parteq}%
\end{equation}
until one reaches self consistency or fixing a value for $\mathbf{x}_{2}$ and
solving for $\mathbf{x}_{1\text{ }}$to obtain%
\begin{equation}
\mathbf{x}_{1}=\mathbb{A}_{11}^{-1}\left(  \mathbf{y}_{1}-\mathbb{A}%
_{12}\mathbf{x}_{2}\right)  \label{x1parteq}%
\end{equation}
until one reaches self consistency. When we have a partial self consistent
solution for $\mathbf{x}_{1}$ and use a fixed $\mathbf{x}_{2}$ we can
substitute back into the equations to obtain
\begin{align}
&  \left(
\begin{array}
[c]{cc}%
\mathbb{A}_{11} & \mathbb{A}_{12}\\
\mathbb{A}_{21} & \mathbb{A}_{22}%
\end{array}
\right)  \left(
\begin{array}
[c]{c}%
\mathbb{A}_{11}^{-1}\left(  \mathbf{y}_{1}-\mathbb{A}_{12}\mathbf{x}%
_{2}\right) \\
\mathbf{x}_{2}%
\end{array}
\right)  -\left(
\begin{array}
[c]{c}%
\mathbf{y}_{1}\\
\mathbf{y}_{2}%
\end{array}
\right) \nonumber\\
&  =\left(
\begin{array}
[c]{c}%
\mathbb{A}_{11}\mathbb{A}_{11}^{-1}\left(  \mathbf{y}_{1}-\mathbb{A}%
_{12}\mathbf{x}_{2}\right)  +\mathbb{A}_{12}\mathbf{x}_{2}\\
\mathbb{A}_{21}\mathbb{A}_{11}^{-1}\left(  \mathbf{y}_{1}-\mathbb{A}%
_{12}\mathbf{x}_{2}\right)  +\mathbb{A}_{22}\mathbf{x}_{2}%
\end{array}
\right)  -\left(
\begin{array}
[c]{c}%
\mathbf{y}_{1}\\
\mathbf{y}_{2}%
\end{array}
\right) \nonumber\\
&  =\left(
\begin{array}
[c]{c}%
\left(  \mathbf{y}_{1}-\mathbb{A}_{12}\mathbf{x}_{2}\right)  +\mathbb{A}%
_{12}\mathbf{x}_{2}\\
\mathbb{A}_{21}\mathbb{A}_{11}^{-1}\left(  \mathbf{y}_{1}-\mathbb{A}%
_{12}\mathbf{x}_{2}\right)  +\mathbb{A}_{22}\mathbf{x}_{2}%
\end{array}
\right)  -\left(
\begin{array}
[c]{c}%
\mathbf{y}_{1}\\
\mathbf{y}_{2}%
\end{array}
\right) \nonumber\\
&  =\left(
\begin{array}
[c]{c}%
\mathbf{0}\\
\mathbb{A}_{21}\mathbb{A}_{11}^{-1}\left(  \mathbf{y}_{1}-\mathbb{A}%
_{12}\mathbf{x}_{2}\right)  +\mathbb{A}_{22}\mathbf{x}_{2}-\mathbf{y}_{2}%
\end{array}
\right)  \label{baksub1}%
\end{align}
While when we have a partial self consistent solution for $\mathbf{x}_{2}$ and
use a fixed $\mathbf{x}_{1}$ we can substitute back into the equations to
obtain
\begin{align}
&  \left(
\begin{array}
[c]{cc}%
\mathbb{A}_{11} & \mathbb{A}_{12}\\
\mathbb{A}_{21} & \mathbb{A}_{22}%
\end{array}
\right)  \left(
\begin{array}
[c]{c}%
\mathbf{x}_{1}\\
\left(  \mathbb{A}_{22}+\mathbb{A}_{21}\mathbb{A}_{11}^{-1}\mathbb{A}%
_{12}\right)  ^{-1}\left(  \mathbf{y}_{2}-\mathbb{A}_{21}\mathbb{A}_{11}%
^{-1}\mathbf{y}_{1}\right)
\end{array}
\right)  -\left(
\begin{array}
[c]{c}%
\mathbf{y}_{1}\\
\mathbf{y}_{2}%
\end{array}
\right) \nonumber\\
&  =\left(
\begin{array}
[c]{c}%
\mathbb{A}_{11}\mathbf{x}_{1}+\mathbb{A}_{12}\left(  \mathbb{A}_{22}%
+\mathbb{A}_{21}\mathbb{A}_{11}^{-1}\mathbb{A}_{12}\right)  ^{-1}\left(
\mathbf{y}_{2}-\mathbb{A}_{21}\mathbb{A}_{11}^{-1}\mathbf{y}_{1}\right) \\
\mathbb{A}_{21}\mathbf{x}_{1}+\mathbb{A}_{22}\left(  \mathbb{A}_{22}%
+\mathbb{A}_{21}\mathbb{A}_{11}^{-1}\mathbb{A}_{12}\right)  ^{-1}\left(
\mathbf{y}_{2}-\mathbb{A}_{21}\mathbb{A}_{11}^{-1}\mathbf{y}_{1}\right)
\end{array}
\right)  -\left(
\begin{array}
[c]{c}%
\mathbf{y}_{1}\\
\mathbf{y}_{2}%
\end{array}
\right) \nonumber\\
&  =\left(
\begin{array}
[c]{c}%
\mathbb{A}_{11}\mathbf{x}_{1}+\mathbb{A}_{12}\mathbb{A}_{22}^{-1}\left\{
\mathbf{y}_{2}-\mathbb{A}_{21}\mathbf{x}_{1}\right\} \\
\mathbf{y}_{2}%
\end{array}
\right)  -\left(
\begin{array}
[c]{c}%
\mathbf{y}_{1}\\
\mathbf{y}_{2}%
\end{array}
\right) \nonumber\\
&  =\left(
\begin{array}
[c]{c}%
\mathbb{A}_{11}\mathbf{x}_{1}+\mathbb{A}_{12}\mathbb{A}_{22}^{-1}\left\{
\mathbf{y}_{2}-\mathbb{A}_{21}\mathbf{x}_{1}\right\}  -\mathbf{y}_{1}\\
\mathbf{0}%
\end{array}
\right)  \label{baksub2}%
\end{align}
The partitioning displayed above obviously only works when the various
inverses exist.

We can obtain equivalent equations by multiplying both sides of eq. $\left(
\ref{ME1}\right)  $ by an invertible matrix e.g.
\begin{equation}
\left(
\begin{array}
[c]{cc}%
\mathbb{B}_{11} & \mathbb{B}_{12}\\
\mathbb{B}_{21} & \mathbb{B}_{22}%
\end{array}
\right)  \left(
\begin{array}
[c]{cc}%
\mathbb{A}_{11} & \mathbb{A}_{12}\\
\mathbb{A}_{21} & \mathbb{A}_{22}%
\end{array}
\right)  \left(
\begin{array}
[c]{c}%
\mathbf{x}_{1}\\
\mathbf{x}_{2}%
\end{array}
\right)  =\left(
\begin{array}
[c]{cc}%
\mathbb{B}_{11} & \mathbb{B}_{12}\\
\mathbb{B}_{21} & \mathbb{B}_{22}%
\end{array}
\right)  \left(
\begin{array}
[c]{c}%
\mathbf{y}_{1}\\
\mathbf{y}_{2}%
\end{array}
\right)  =\left(
\begin{array}
[c]{c}%
\mathbf{z}_{1}\\
\mathbf{z}_{2}%
\end{array}
\right)  \label{pre}%
\end{equation}
and then partitioning in the same manner. If $\mathbb{A}_{12}=\mathbb{A}%
_{21}=0$ then the equations are uncoupled if they are linear, but are still
coupled if they are nonlinear. If $\mathbb{A}_{11}=\mathbb{A}_{22}=0$ then one
can premultiply eq. $\left(  \ref{pre}\right)  $ by $\left(
\begin{array}
[c]{cc}%
\mathbb{O} & \mathbb{I}\\
\mathbb{I} & \mathbb{O}%
\end{array}
\right)  $ and then solve, but then the partial solutions do not give the same
back substituted properties displayed in eqs. $\left(  \ref{baksub1}\right)  $
and $\left(  \ref{baksub2}\right)  $ i.e. the residual vectors of the original
problem generally have zero components in different positions or no zero
components at all.

\subsection{Zero $\mathbb{A}_{22}$ Block}

The \emph{nonlinear }equation eq. $\left(  \ref{ME1}\right)  $ for
$\mathbf{x}_{1}$ and $\mathbf{x}_{2}$
\begin{equation}
\left(
\begin{array}
[c]{cc}%
\mathbb{A}_{11} & \mathbb{A}_{12}\\
\mathbb{A}_{21} & \mathbb{A}_{22}%
\end{array}
\right)  \left(
\begin{array}
[c]{c}%
\mathbf{x}_{1}\\
\mathbf{x}_{2}%
\end{array}
\right)  -\left(
\begin{array}
[c]{c}%
\mathbf{y}_{1}\\
\mathbf{y}_{2}%
\end{array}
\right)  =0
\end{equation}
does not have be solved in the preceding manner, in fact when one of the
diagonal blocks is zero it may become impossible. In particular consider
\begin{equation}
\left(
\begin{array}
[c]{cc}%
0 & \mathbb{A}_{12}\\
\mathbb{A}_{21} & \mathbb{A}_{22}%
\end{array}
\right)  \left(
\begin{array}
[c]{c}%
\mathbf{x}_{1}\\
\mathbf{x}_{2}%
\end{array}
\right)  -\left(
\begin{array}
[c]{c}%
\mathbf{y}_{1}\\
\mathbf{y}_{2}%
\end{array}
\right)  =0\Longleftrightarrow\left.
\begin{array}
[c]{c}%
\qquad\qquad\mathbb{A}_{12}\mathbf{x}_{2}-\mathbf{y}_{1}=0\\
\mathbb{A}_{21}\mathbf{x}_{1}+\mathbb{A}_{22}\mathbf{x}_{2}-\mathbf{y}_{2}=0
\end{array}
\right\}
\end{equation}
which gives for fixed $\mathbf{x}_{1}$ the equation
\begin{equation}
\mathbb{A}_{22}\mathbf{x}_{2}=\mathbf{y}_{2}-\mathbb{A}_{21}\mathbf{x}_{1}
\label{x2le}%
\end{equation}
for $\mathbf{x}_{2}$. A self consistent solution, $\mathbf{x}_{2\#}\left(
\mathbf{x}_{1\left(  k\right)  }\right)  $, is obtained by iterating
$\mathbf{x}_{2}$ in eq. $\left(  \ref{x2le}\right)  $ at the fixed value
$\mathbf{x}_{1}=\mathbf{x}_{1\left(  k\right)  }$
\begin{equation}
\mathbf{x}_{2\#}\left(  \mathbf{x}_{1\left(  k\right)  }\right)
=\mathbb{A}_{22}^{-1}\left(  \mathbf{y}_{2}-\mathbb{A}_{21}\mathbf{x}%
_{1\left(  k\right)  }\right)  \label{x2}%
\end{equation}
If one substitutes back this partial solution, noting that all quantities,
$\mathbb{A}_{12},\mathbb{A}_{21},\mathbb{A}_{22},\mathbf{y}_{1}$ and
$\mathbf{y}_{2}$ are evaluated at $\left(  \mathbf{x}_{1\left(  k\right)
},\mathbf{x}_{2\#}\left(  \mathbf{x}_{1\left(  k\right)  }\right)  \right)  $,
we obtain
\begin{align}
&  \left(
\begin{array}
[c]{cc}%
0 & \mathbb{A}_{12}\\
\mathbb{A}_{21} & \mathbb{A}_{22}%
\end{array}
\right)  \left(
\begin{array}
[c]{c}%
\mathbf{x}_{1\left(  k\right)  }\\
\mathbb{A}_{22}^{-1}\left(  \mathbf{y}_{2}-\mathbb{A}_{21}\mathbf{x}_{1\left(
k\right)  }\right)
\end{array}
\right)  -\left(
\begin{array}
[c]{c}%
\mathbf{y}_{1}\\
\mathbf{y}_{2}%
\end{array}
\right)  =\nonumber\\
&  \qquad\left(
\begin{array}
[c]{c}%
\mathbb{A}_{12}\mathbb{A}_{22}^{-1}\left(  \mathbf{y}_{2}-\mathbb{A}%
_{21}\mathbf{x}_{1\left(  k\right)  }\right)  -\mathbf{y}_{1}\\
\mathbb{A}_{21}\mathbf{x}_{1\left(  k\right)  }+\mathbf{y}_{2}-\mathbb{A}%
_{21}\mathbf{x}_{1\left(  k\right)  }-\mathbf{y}_{2}%
\end{array}
\right)  =\left(
\begin{array}
[c]{c}%
\mathbb{A}_{12}\mathbb{A}_{22}^{-1}\left(  \mathbf{y}_{2}-\mathbb{A}%
_{21}\mathbf{x}_{1\left(  k\right)  }\right)  -\mathbf{y}_{1}\\
\mathbf{0}%
\end{array}
\right)  \label{baksub}%
\end{align}
The self consistent solution, $\mathbf{x}_{1\#\left(  k+1\right)  }$, for
$\mathbf{x}_{1}$ at $\mathbf{x}_{2\#}\left(  \mathbf{x}_{1\left(  k\right)
}\right)  $ is then obtained by iterating $\mathbf{x}_{1\left(  k+1\right)  }$
in%
\begin{equation}
\mathbb{A}_{12}\mathbb{A}_{22}^{-1}\left(  \mathbf{y}_{2}-\mathbb{A}%
_{21}\mathbf{x}_{1\left(  k+1\right)  }\right)  -\mathbf{y}_{1}=0
\end{equation}
or in the formal relationship
\begin{equation}
\mathbf{x}_{1\left(  k+1\right)  }=\mathbb{A}_{12}^{-1}\mathbb{A}%
_{22}\mathbb{A}_{21}^{-1}\left(  \mathbf{y}_{1}-\mathbb{A}_{12}\mathbb{A}%
_{22}^{-1}\mathbf{y}_{2}\right)  \label{x1kplus1}%
\end{equation}
where all quantities are evaluated at $\left(  \mathbf{x}_{1\left(
k+1\right)  },\mathbf{x}_{2\#}\left(  \mathbf{x}_{1\left(  k\right)  }\right)
\right)  $ and left inverses are used for $\mathbb{A}_{12}$ and $\mathbb{A}%
_{21}$. These two coupled iterative sequences can then be repeated until self
consistency is reached for both $\mathbf{x}_{1}$ and $\mathbf{x}_{2}$ to
obtain a solution of the full non linear problem.

If $\mathbb{A}_{22}$ is singular then eq. $\left(  \ref{x2}\right)  $ becomes%
\begin{equation}
\mathbf{x}_{2\#}\left(  \mathbf{x}_{1\left(  k\right)  },\mathbf{k}\right)
=\mathbb{A}_{22}^{-1}\left(  \mathbf{y}_{2}-\mathbb{A}_{21}\mathbf{x}%
_{1\left(  k\right)  }\right)  +\mathbf{k} \label{x2eq}%
\end{equation}
where we use the generalized inverse $\mathbb{A}_{22}^{-1}$ of $\mathbb{A}%
_{22}$, $\mathbf{k}\in\ker\mathbb{A}_{22}$ and the subsidiary condition
\begin{equation}
P_{\ker\mathbb{A}_{22}}\left(  \mathbf{y}_{2}-\mathbb{A}_{21}\mathbf{x}%
_{1\left(  k\right)  }\right)  =0 \label{x2cond}%
\end{equation}
must be satisfied, with all terms evaluated at $\left(  \mathbf{x}_{1\left(
k\right)  },\mathbf{x}_{2\#}\left(  \mathbf{x}_{1\left(  k\right)
},\mathbf{k}\right)  \right)  $. Then the back substitution eq. $\left(
\ref{baksub}\right)  $ becomes
\begin{align}
&  \left(
\begin{array}
[c]{cc}%
0 & \mathbb{A}_{12}\\
\mathbb{A}_{21} & \mathbb{A}_{22}%
\end{array}
\right)  \left(
\begin{array}
[c]{c}%
\mathbf{x}_{1\left(  k\right)  }\\
\mathbb{A}_{22}^{-1}\left(  \mathbf{y}_{2}-\mathbb{A}_{21}\mathbf{x}_{1\left(
k\right)  }\right)  +\mathbf{k}%
\end{array}
\right)  -\left(
\begin{array}
[c]{c}%
\mathbf{y}_{1}\\
\mathbf{y}_{2}%
\end{array}
\right)  =\nonumber\\
&  \qquad\left(
\begin{array}
[c]{c}%
\mathbb{A}_{12}\mathbb{A}_{22}^{-1}\left(  \mathbf{y}_{2}-\mathbb{A}%
_{21}\mathbf{x}_{1\left(  k\right)  }\right)  +\mathbb{A}_{12}\mathbf{k}%
-\mathbf{y}_{1}\\
P_{\ker\mathbb{A}_{22}}^{c}\left(  \mathbf{y}_{2}-\mathbb{A}_{21}%
\mathbf{x}_{1\left(  k\right)  }\right)  +\mathbb{A}_{21}\mathbf{x}_{1\left(
k\right)  }-\mathbf{y}_{2}%
\end{array}
\right)  =\nonumber\\
&  \qquad\qquad\qquad\qquad\left(
\begin{array}
[c]{c}%
\mathbb{A}_{12}\mathbb{A}_{22}^{-1}\left(  \mathbf{y}_{2}-\mathbb{A}%
_{21}\mathbf{x}_{1\left(  k\right)  }\right)  +\mathbb{A}_{12}\mathbf{k}%
-\mathbf{y}_{1}\\
\mathbf{0}%
\end{array}
\right)  \label{subbak}%
\end{align}
for any $\mathbf{k}$. The equation for a self consistent value, $\mathbf{x}%
_{1\#\left(  k+1\right)  }\left(  \mathbf{x}_{2\#}\left(  \mathbf{x}_{1\left(
k\right)  },\mathbf{k}\right)  \right)  $, for $\mathbf{x}_{1}$ at
$\mathbf{x}_{2}=\mathbf{x}_{2\#}\left(  \mathbf{x}_{1\left(  k\right)
},\mathbf{k}\right)  $ obtained from eq. $\left(  \ref{subbak}\right)  $ is%
\begin{equation}
\mathbb{A}_{12}\mathbb{A}_{22}^{-1}\mathbb{A}_{21}\mathbf{x}_{1\#\left(
k+1\right)  }=\mathbb{A}_{12}\mathbb{A}_{22}^{-1}\mathbf{y}_{2}-\mathbf{y}%
_{1}+\mathbb{A}_{12}\mathbf{k} \label{scx1}%
\end{equation}
which can be formally solved as
\begin{equation}
\mathbf{x}_{1\#\left(  k+1\right)  }\left(  \mathbf{x}_{1\left(  k\right)
},\mathbf{k}\right)  =\mathbb{A}_{21}^{-1}\mathbb{A}_{22}\mathbb{A}_{12}%
^{-1}\left(  \mathbb{A}_{12}\mathbb{A}_{22}^{-1}\mathbf{y}_{2}-\mathbf{y}%
_{1}\right)  =\mathbb{A}_{21}^{-1}\mathbf{y}_{2}-\mathbb{A}_{21}%
^{-1}\mathbb{A}_{22}\mathbb{A}_{12}^{-1}\mathbf{y}_{1}%
\end{equation}
where $\mathbb{A}_{12},\mathbb{A}_{21},\mathbb{A}_{22},\mathbf{y}_{1},$ and
$\mathbf{y}_{2}$ are evaluated at $\left(  \mathbf{x}_{1\#\left(  k+1\right)
},\mathbf{x}_{2\#}\left(  \mathbf{x}_{1\left(  k\right)  },\mathbf{k}\right)
\right)  $. This value can then be substituted back into the equation for
$\mathbf{x}_{2}$, which can then be iterated to self consistency once more,
the new self consistent value of $\mathbf{x}_{2}$ can then be used to obtain a
new self consistent value of $\mathbf{x}_{1}$ and so on.

Instead\ of iterating to self consistency for a new\ $\mathbf{x}_{1\#\left(
k+1\right)  }$ and then substituting back into the equation for $\mathbf{x}%
_{2}$, which would be the most efficient way to solve the total coupled
problem, one could find paths of self consistent solutions, $\mathbf{x}%
_{2\#}\left(  \mathbf{x}_{1},\mathbf{k}\right)  $, for all possible
$\mathbf{x}_{1}$'s producing paths, $\mathbf{x}_{2\#}\left(  \mathbf{x}%
_{1},\mathbf{k}\right)  $, of $\mathbf{x}_{2}$'s that satisfy
\begin{equation}
\mathbf{x}_{2\#}\left(  \mathbf{x}_{1},\mathbf{k}\right)  =\mathbb{A}%
_{22}^{-1}\left(  \mathbf{y}_{2}-\mathbb{A}_{21}\mathbf{x}_{1}\right)
+\mathbf{k}%
\end{equation}
where
\begin{align}
\mathbb{A}_{22}  &  =\mathbb{A}_{22}\left(  \mathbf{x}_{1},\mathbf{x}%
_{2\#}\left(  \mathbf{x}_{1},\mathbf{k}\right)  \right) \nonumber\\
\mathbf{y}_{2}  &  =\mathbf{y}_{2}\left(  \mathbf{x}_{1},\mathbf{x}%
_{2\#}\left(  \mathbf{x}_{1},\mathbf{k}\right)  \right) \nonumber\\
\mathbb{A}_{21}  &  =\mathbb{A}_{21}\left(  \mathbf{x}_{1},\mathbf{x}%
_{2\#}\left(  \mathbf{x}_{1},\mathbf{k}\right)  \right) \nonumber\\
\mathbf{k}  &  =\mathbf{k}\left(  \mathbf{x}_{1}\right)
\end{align}
and
\begin{align}
\left(
\begin{array}
[c]{cc}%
0 & \mathbb{A}_{12}\\
\mathbb{A}_{21} & \mathbb{A}_{22}%
\end{array}
\right)  \left(
\begin{array}
[c]{c}%
\mathbf{x}_{1}\\
\mathbf{x}_{2\#}\left(  \mathbf{x}_{1},\mathbf{k}\right)
\end{array}
\right)  -\left(
\begin{array}
[c]{c}%
\mathbf{y}_{1}\\
\mathbf{y}_{2}%
\end{array}
\right)   &  =\left(
\begin{array}
[c]{c}%
\mathbb{A}_{12}\mathbf{x}_{2\#}\left(  \mathbf{x}_{1},\mathbf{k}\right)
-\mathbf{y}_{1}\\
\mathbb{A}_{21}\mathbf{x}_{1}+\mathbb{A}_{22}\mathbf{x}_{2\#}\left(
\mathbf{x}_{1},\mathbf{k}\right)  -\mathbf{y}_{2}%
\end{array}
\right) \nonumber\\
&  =\left(
\begin{array}
[c]{c}%
\mathbb{A}_{12}\mathbf{x}_{2\#}\left(  \mathbf{x}_{1},\mathbf{k}\right)
-\mathbf{y}_{1}\\
0
\end{array}
\right)
\end{align}
where
\begin{align}
\mathbb{A}_{12}  &  =\mathbb{A}_{12}\left(  \mathbf{x}_{1},\mathbf{x}%
_{2\#}\left(  \mathbf{x}_{1},\mathbf{k}\right)  \right) \nonumber\\
\mathbf{y}_{1}  &  =\mathbf{y}_{1}\left(  \mathbf{x}_{1},\mathbf{x}%
_{2\#}\left(  \mathbf{x}_{1},\mathbf{k}\right)  \right)
\end{align}
Then solve the nonlinear equation for $\mathbf{x}_{1}$
\begin{align}
&  \mathbb{A}_{12}\mathbf{x}_{2\#}\left(  \mathbf{x}_{1},\mathbf{k}\right)
-\mathbf{y}_{1}=0\nonumber\\
&  \equiv\mathbb{A}_{12}\mathbb{A}_{22}^{-1}\left(  \mathbf{y}_{2}%
-\mathbb{A}_{21}\mathbf{x}_{1}\right)  +\mathbb{A}_{12}\mathbf{k}%
-\mathbf{y}_{1}=0 \label{X1eq}%
\end{align}
with formal solution
\[
\mathbf{x}_{1\#}=\mathbb{A}_{21}^{-1}\mathbf{y}_{2}-\mathbb{A}_{21}%
^{-1}\mathbb{A}_{22}\mathbb{A}_{12}^{-1}\mathbf{y}_{1}=\mathbb{A}_{21}%
^{-1}\left(  \mathbf{y}_{2}-\mathbb{A}_{22}\mathbb{A}_{12}^{-1}\mathbf{y}%
_{1}\right)
\]
(which corresponds to the equation we identified for $\mathcal{V}_{den}$ in
eq. $\left(  \ref{gammaeq}\right)  $) to determine the exact $\mathbf{x}_{1}$,
where everything in eq. $\left(  \ref{X1eq}\right)  $ is evaluated along the
paths $\left(  \mathbf{x}_{1},\mathbf{x}_{2\#}\left(  \mathbf{x}%
_{1},\mathbf{k}\right)  \right)  $. The different paths of solutions
$\mathbf{x}_{2\#}\left(  \mathbf{x}_{1},\mathbf{k}\right)  $ are labelled by
$\mathbf{k}$ i.e.
\begin{equation}
\mathbf{x}_{2\#}\left(  \mathbf{x}_{1},\mathbf{k}\right)  \neq\mathbf{x}%
_{2\#}\left(  \mathbf{x}_{1},\mathbf{k}^{\prime}\right)  \text{ for
}\mathbf{k}\left(  \mathbf{x}_{1}\right)  \neq\mathbf{k}^{\prime}\left(
\mathbf{x}_{1}\right)
\end{equation}
and an extra criterion must be used to determine unique paths. As
$\mathbf{x}_{2\#}\left(  \mathbf{x}_{1},\mathbf{k}\right)  $ is always a self
consistent solution for $\mathbf{x}_{2}$ in eq. $\left(  \ref{x2eq}\right)  $
one can observe that for any $\mathbf{x}_{1}$ we have
\begin{align}
&  \left(
\begin{array}
[c]{cc}%
0 & \mathbb{A}_{12}\\
\mathbb{A}_{21} & \mathbb{A}_{22}%
\end{array}
\right)  \left(
\begin{array}
[c]{c}%
\mathbf{x}_{1}\\
\mathbf{x}_{2\#}\left(  \mathbf{x}_{1},\mathbf{k}\right)
\end{array}
\right)  -\left(
\begin{array}
[c]{c}%
\mathbf{y}_{1}\\
\mathbf{y}_{2}%
\end{array}
\right)  =\left(
\begin{array}
[c]{c}%
\mathbb{A}_{12}\mathbf{x}_{2\#}\left(  \mathbf{x}_{1}\right)  -\mathbf{y}%
_{1}\\
\mathbb{A}_{21}\mathbf{x}_{1}+\mathbb{A}_{22}\mathbf{x}_{2\#}\left(
\mathbf{x}_{1},\mathbf{k}\right)  -\mathbf{y}_{2}%
\end{array}
\right) \nonumber\\
&  =\left(
\begin{array}
[c]{c}%
\mathbb{A}_{12}\mathbf{x}_{2\#}\left(  \mathbf{x}_{1}\right)  -\mathbf{y}%
_{1}\\
\mathbb{A}_{21}\mathbf{x}_{1}+\mathbb{A}_{22}\left(  \mathbb{A}_{22}%
^{-1}\left(  \mathbf{y}_{2}-\mathbb{A}_{21}\mathbf{x}_{1}\right)
+\mathbf{k}\right)  -\mathbf{y}_{2}%
\end{array}
\right)  =\left(
\begin{array}
[c]{c}%
\mathbb{A}_{12}\mathbf{x}_{2\#}\left(  \mathbf{x}_{1}\right)  -\mathbf{y}%
_{1}\\
\mathbf{0}%
\end{array}
\right)
\end{align}
i.e.%
\begin{equation}
\left(
\begin{array}
[c]{cc}%
0 & \mathbb{A}_{12}\\
\mathbb{A}_{21} & \mathbb{A}_{22}%
\end{array}
\right)  \left(
\begin{array}
[c]{c}%
\mathbf{x}_{1}\\
\mathbf{x}_{2\#}\left(  \mathbf{x}_{1}\right)
\end{array}
\right)  -\left(
\begin{array}
[c]{c}%
\mathbf{y}_{1}\\
\mathbf{y}_{2}%
\end{array}
\right)  =\left(
\begin{array}
[c]{c}%
\mathbb{A}_{12}\mathbf{x}_{2\#}\left(  \mathbf{x}_{1}\right)  -\mathbf{y}%
_{1}\\
0
\end{array}
\right)
\end{equation}



\end{document}